\chapter{Background} \label{ch:background}



\section{Explainability} \label{sec:explainability}

Explainability is the field that studies and designs methods to make a system understandable to humans, intended as a spectrum of different stakeholders with different goals and needs. In fact, across different domains, explanations can be used to improve model performance, enhance and support human decision-making, or analyze biases \cite{2024_xai_manifesto}. It is a multidisciplinary field that spans across different areas such as computer science, psychology, and ethics \cite{2020_xai_survey}. Indeed, explainability not only has to be technically functional and correct, but should also account for users' needs and rights \cite{2021_xai_multidisciplinary}. Moreover, the necessity to provide an explanation for the output of a model is also mentioned in some legislations. For instance, in the \ac{eu}, explainability is mentioned as a right in article 22 of the \ac{gdpr} \cite{gdpr_22} and it is one of the principles of trustworthy \ac{ai} \cite{2019_eu_trustworthy_ai}.

Different taxonomies can be used to describe explainability methods. From a high-level classification, an explainability method can be employed at the pre-modelling, \antehoc modelling, or \posthoc modelling levels \cite{2021_mm_xai_survey, 2024_mm_xai_survey, 2024_xai_manifesto}. Pre-modelling explainability involves methods used before developing the model and it is mainly used to explore and analyze the data through dimensionality reduction or feature selection techniques. \Antehoc explainability aims at designing models that are intrinsically explainable; this includes models such as $k$-nearest neighbors, decision trees, and other more advanced approaches. \Posthoc explainability, which this thesis focuses on, includes methods that are applied on a model after development; this includes a variety of methods such as perturbation, gradient-based, or proxy methods that will be further discussed in the following sections \cite{2021_mm_xai_survey, 2024_mm_xai_survey}.
More in-depth, an explanation can be local, if it works on single instances, or global, if it aims at explaining the whole model \cite{2021_mm_xai_survey}.
Furthermore, in the case of \posthoc explainability, a method can be model-specific or model-agnostic depending on whether it is independent of the underlying model or not \cite{2021_mm_xai_survey, 2024_xai_manifesto}.

In practical terms, this work mainly focuses on \posthoc local explainability for neural networks. In particular, the proposed method falls within the class of attribution-based methods, which aim at computing importance scores for each input feature with respect to a specific model prediction. These methods produce explanations in the form of relevance maps or saliency scores, indicating how much each feature contributed to the output for a given instance. 
Several existing approaches have been proposed in the literature. In this section, we briefly describe them and provide more details in \Cref{sec:baselines} when describing our baselines. One of the simplest methods is based on Saliency \cite{saliency}, which builds on the idea of computing the gradient of the model output with respect to the input features and interpret its magnitude as a measure of sensitivity of how small changes in the input space affect the prediction. Although conceptually simple and computationally efficient, a major drawback of this method is that it often suffers from noise and can produce explanations that are unstable or difficult to interpret.
Another method is DeepLIFT \cite{deeplift}, which propagates attribution scores backward through the network by comparing the activation of each neuron to a reference activation. Compared to simpler methods, DeepLIFT provides more stability and avoids some gradient instability issues. Building on top of it, DeepLIFT-SHAP \cite{deepliftshap} uses DeepLIFT to approximate Shapley values \cite{shapley_values}, which are a game theory notion based on the idea of fairly distributing feature importance, but has the drawback of being computationally expensive to compute exactly. A drawback of DeepLIFT and other methods based on it is that it is sensitive to the choice of baselines that are used as reference activations. Therefore, in some cases, it can produce drastically different results depending on which prior features are provided.
Integrated Gradients \cite{lig} is another well-established method that approximates the integral of the gradient along a path from some given baseline inputs to the actual input, satisfying desirable axioms such as sensitivity and implementation invariance. Again, the choice of baselines is a critical choice that can significantly influence the resulting explanations.
A closely related approach is Gradient-SHAP \cite{gradient_shap}, which combines ideas from Integrated Gradients with stochastic sampling. Instead of integrating along a single path from a baseline to the input, Gradient-SHAP samples random interpolation points between multiple baseline inputs and the actual input. By averaging gradients over these paths, the method approximates Shapley values while reducing variance compared to single-baseline techniques, therefore improving on robustness but with an increase in computational cost due to the required sampling.
Finally, Guided Backpropagation \cite{guided_backprop} modifies the standard backpropagation procedure by suppressing negative gradients during the backward pass through ReLU activations. Concretely, only gradients associated with positive forward activations and positive backward signals are propagated. This heuristic approach often produces sharp and visually appealing explanations, particularly in image-based settings. However, it does not necessarily provide faithful feature importance scores and does not have theoretical grounding.


\section{Motivation}

In \Cref{sec:explainability}, we illustrated the most commonly known attribution-based explainability methods in the literature. However, each of these methods has some issues such as sensitivity to the choice of parameters and to the choice of prior baselines. Another possible drawback is that some of these methods require multiple iterations or are based on computing the gradient by performing backpropagation along the whole length of the model at inference time to produce the attribution map, which can make the applicability of these approaches slow in practical terms.

Considering these observations, in this thesis, we propose a method based on learning a dedicated neural network to produce attribution maps, which aims at reducing some of the issues we mentioned. Firstly, we remove the need to provide some arbitrary baselines. Secondly, we remove the need of choosing method-specific hyperparameters and instead reduce it to the choice of more traditional and familiar hyperparameters for training a neural network. Thirdly, we remove some of the computational cost at inference time as a single forward pass of a smaller model produces the attribution map. Lastly, although not strictly related to the other methods we mentioned, our approach does not require having ground-truth explanations as it is trained in a self-supervised way, so that it can be applied to a wide range of scenarios. 

Obviously, the approach we are proposing has its own drawbacks: it moves most of the computational cost from inference to training time and, instead of providing a limited number of baseline inputs, it requires the whole, or a large subset, of original training data of the model which is usually larger in size. Nevertheless, we argue that these limitations are negligible and still make this approach practically feasible and convenient as it does not require collecting new data and the dedicated model to produce attribution maps does not have to be a large model.



\section{Related Work}

In the literature, there are some approaches closely related to the idea we presented in the previous section.
In the context of large language models, \citeauthor{barkan-etal-2024-llm} \cite{barkan-etal-2024-llm} propose a framework called Attributive Masking Learning where they train an auxiliary model to mask input tokens while maintaining the output of the model as close as possible to the original one.
Instead, \citeauthor{bhattacharya2024gradientfreeposthocexplainabilityusing} \cite{bhattacharya2024gradientfreeposthocexplainabilityusing} propose a distillation approach, where a dedicated model produces attribution maps that are then used to mask the input and fed through a student network that has to mimic the original model to explain.
\citeauthor{liu2019explainablenlpgenerativeexplanation} \cite{liu2019explainablenlpgenerativeexplanation} also build on the same idea of learning a dedicated explanation model in the case of text classification. However, their approach requires knowing the ground-truth explanations, which is not always possible to have.
Finally, \citeauthor{kanehira2019learningexplaincomplementalexamples} \cite{kanehira2019learningexplaincomplementalexamples} propose a framework in the context of image classification that consists of learning two dedicated networks to provide a textual explanation and visual examples, respectively. However, it also requires ground-truth information, which makes it hard to apply in practice.

Overall, to the best of our knowledge, there is no prior work in the literature that proposes a data-agnostic learned explainability method for producing attribution maps that works across different data types and does not require having ground-truth of the explanations.
